\begin{table*}[t]
\centering
\caption{Comparison with four traditional planning baselines on the held-out
TC segments from the same 500-TC train/test split. Surface values are means
over the 4/8/12/16-SUV fleets, both speed limits, and all five deployment
radii. Subsurface values are means over all seven depths and both propulsion
limits. Smaller errors and paths are better; higher current and coverage are
better. No single method is best in all subsurface columns.}
\label{tab:classical_planners}
\small
\begin{tabular*}{\textwidth}{@{\extracolsep{\fill}}lrrr@{}}
\toprule
\multicolumn{4}{l}{\textbf{Panel A: surface tracking}}\\
Method & Final error (km) & Time-averaged distance (km) & Path per SUV (km)\\
\midrule
Parameter-sharing IPPO & \textbf{341.7} & \textbf{349.1} & \textbf{463.6} \\
Potential field & 416.5 & 405.9 & 727.6 \\
Receding-horizon A* & 459.0 & 430.4 & 824.9 \\
Lawnmower coverage & 554.8 & 478.4 & 508.4 \\
Information greedy & 467.0 & 429.6 & 525.2 \\
\midrule
\multicolumn{4}{l}{\textbf{Panel B: subsurface observation}}\\
Method & Mean current (m/s) & Nominal coverage (km$^2$) & Path (km)\\
\midrule
Parameter-sharing IPPO & 0.409 & \textbf{3630} & 126.4 \\
Potential field & \textbf{0.449} & 1976 & 112.5 \\
Receding-horizon A* & 0.428 & 2275 & \textbf{94.8} \\
Lawnmower coverage & 0.409 & 2902 & 111.1 \\
Information greedy & 0.432 & 1520 & 102.4 \\
\bottomrule
\end{tabular*}
\end{table*}
